% !Mode:: "TeX:UTF-8"
\chapter{总结与展望}
\label{cha:summary}

\section{研究工作总结}

本文围绕点云补全任务中的核心问题展开研究，针对残缺点云补全过程中出现的密度失衡与分布不均匀问题，提出了一套基于Transformer的点云补全算法优化方案，并构建了从补全点云到三维网格模型的完整表面重建流程。本文的主要研究工作与贡献总结如下：

（1）提出了基于混合损失函数的点云补全优化策略。在粗粒度生成阶段，引入EMD（Earth Mover's Distance）损失替代单一的标准CD损失，通过双射匹配约束强化了骨架点云的全局拓扑对齐能力。在细粒度优化阶段，设计了包含密度感知倒角距离（DCD）和排斥损失的混合损失函数，通过密度感知权重机制和局部排斥力约束，有效改善了生成点云的分布均匀性。分阶段权重调度策略使得模型能够在不同训练阶段获得最合适的优化目标，实现了多损失之间的协同优化。

（2）引入了自适应查询生成机制。通过自适应查询排序和去噪任务辅助训练两个组件，提升了查询点的质量和模型对查询点特征的利用效率。查询排序机制通过学习性重要性评估，使得模型能够自动识别几何特征丰富的关键区域；去噪任务则通过额外的监督信号增强了模型对输入扰动的容忍度，提升了泛化性能。

（3）构建了完整的三维表面重建流程。基于泊松表面重建算法，实现了从补全点云到三角网格模型的转换。该流程涵盖法线估计与定向、泊松重建、密度阈值过滤和Taubin平滑等关键步骤，能够基于补全后的高质量点云生成表面连续、细节丰富的三维网格模型，为点云补全算法的实际应用提供了有效的下游任务验证。

（4）通过系统的实验验证了算法的有效性。在ShapeNet-34数据集上的对比实验表明，本算法在CD-L2和F-Score@1\%指标上均取得了具有竞争力的结果，优于PoinTr、SnowflakeNet、SeedFormer等经典方法。消融实验进一步验证了各项优化措施的独立贡献，三项优化措施累计使CD-L2指标降低35.8\%（从$1.23 \times 10^{-3}$降至$0.79 \times 10^{-3}$），F-Sore@1\%提升8.1\%（从42.1\%增至45.5\%）。实验结果充分证明了本文所提优化策略的有效性和合理性。

\section{未来研究展望}

尽管本文在点云补全算法优化方面取得了一定的研究成果，但受限于研究时间和实验条件，仍存在若干值得进一步探索的方向：

（1）架构层面的进一步优化。从对比实验结果可以看出，本算法与当前SOTA级模型（如AnchorFormer、PCD-MT等）仍存在一定的性能差距。未来工作可以考虑在架构设计方面引入创新性机制，例如探索Mamba等新型状态空间模型与Transformer的融合，或者设计更加高效的特征聚合策略，以进一步提升补全精度和计算效率。

（2）面向复杂场景的泛化能力提升。本文的实验评估主要基于ShapeNet-34和Projected-ShapeNet数据集，这些数据集的残缺模式相对规则。未来工作可以在更加复杂和真实的残缺场景下验证算法的鲁棒性，例如处理由真实深度相机采集的点云数据，或者应对极端遮挡、严重噪声等挑战性条件。此外，探索跨类别的零样本补全能力也是一个有价值的研究方向。

（3）下游应用的拓展。点云补全作为三维视觉领域的基础任务，其下游应用涵盖三维重建、机器人抓取、自动驾驶、虚拟现实等多个方向。未来工作可以针对特定的应用场景，探索点云补全算法与下游任务的联合优化策略，实现端到端的性能提升。例如，在机器人抓取任务中，补全质量直接影响抓取姿态的估计精度；在自动驾驶场景中，补全算法可以帮助构建更加完整的环境三维表示。